%==================================================================
%  CV - Emmanuel Narro  |  Desarrollador de Software Jr.
%  Compilar con: pdflatex CV_Emmanuel_Narro.tex   (o en Overleaf)
%  Basado en el portafolio: github.com/Emmanuel121003
%==================================================================
\documentclass[11pt,a4paper]{article}

\usepackage[utf8]{inputenc}
\usepackage[T1]{fontenc}
\usepackage{textcomp}
\usepackage[spanish,es-noshorthands]{babel}
\usepackage[margin=1.6cm]{geometry}
\usepackage{titlesec}
\usepackage{enumitem}
\usepackage{xcolor}
\usepackage{hyperref}
\usepackage{lmodern}
\usepackage{microtype}

% ---- Paleta (cobre / tinta, a juego con el portafolio) ----
\definecolor{accent}{HTML}{A85A22}
\definecolor{ink}{HTML}{1A2130}
\definecolor{muted}{HTML}{5A6472}

\hypersetup{colorlinks=true, urlcolor=accent, linkcolor=accent}

% ---- Formato de secciones ----
\titleformat{\section}
  {\color{accent}\bfseries\large\scshape}
  {}{0pt}{}[{\color{accent!40}\titlerule[1pt]}]
\titlespacing{\section}{0pt}{10pt}{6pt}

\setlist[itemize]{leftmargin=1.2em, itemsep=1pt, topsep=2pt, parsep=0pt}
\pagestyle{empty}
\setlength{\parindent}{0pt}

% ---- Comando para entradas de experiencia ----
\newcommand{\entry}[4]{%
  \textbf{#1} \hfill {\color{muted}#2}\\[1pt]
  {\color{accent}\itshape #3} \hfill {\color{muted}\itshape #4}\\[2pt]
}

\begin{document}

%==================== ENCABEZADO ====================
\begin{center}
  {\color{ink}\Huge\bfseries Emmanuel Narro}\\[4pt]
  {\large\color{accent} Desarrollador de Software Jr. \textbar{} Backend \& Full Stack}\\[6pt]
  {\color{muted}\small
   Guadalajara, Jalisco, México \enspace\textbullet\enspace
   \href{https://wa.me/526182704889}{+52 618 270 4889 (WhatsApp)} \enspace\textbullet\enspace
   \href{mailto:emmanuel.nare.12@gmail.com}{emmanuel.nare.12@gmail.com}\\[3pt]
   \href{https://www.linkedin.com/in/emmanuel-narro-8b4476170}{linkedin.com/in/emmanuel-narro-8b4476170}
   \enspace\textbullet\enspace
   \href{https://github.com/Emmanuel121003}{github.com/Emmanuel121003}}
\end{center}

%==================== PERFIL ====================
\section{Perfil profesional}
Ingeniero en Desarrollo y Gestión de Software enfocado en \textbf{backend y desarrollo web full stack}, con
experiencia profesional construyendo y dando soporte a \textbf{sistemas empresariales en producción}. Trabajo con
\textbf{.NET/C\#, ASP.NET Web API, Python y SQL Server}, integraciones mediante \textbf{APIs REST} y desarrollo
multiplataforma (web, móvil e IoT). He cubierto el ciclo completo: modelo de datos, lógica de negocio, despliegue en
sitio y resolución de incidencias. Busco un rol Junior de Full Stack, Backend o Software Developer.

%==================== EXPERIENCIA ====================
\section{Experiencia profesional}

\entry{E-GO Desarrollo de Software}{Guadalajara, Jal.}{Desarrollador de Software}{Mayo 2026 -- Actualidad}
\begin{itemize}
  \item \textbf{Scale II} (control de peso industrial, en producción): principal contribuidor de la app Android (Java)
        e integración con la aplicación de escritorio en .NET/C\#, comunicación con hardware (Raspberry Pi) y SQL Server.
  \item \textbf{Gens} (punto de venta y e-commerce): contribuidor del POS de escritorio en C\#/.NET y del frontend
        (Vue.js) y servicios web (Node.js) de la plataforma web.
  \item \textbf{PVM} (punto de venta móvil): desarrollo de la plataforma web en PHP y aporte a la app Android.
  \item Despliegue en sitio y atención de incidencias; sincronización y respaldo de datos vía APIs REST y
        optimización de consultas SQL.
  \item Trabajo en equipo con Scrum, Git y GitHub; documentación técnica.
\end{itemize}

\entry{E-GO Desarrollo de Software}{Guadalajara, Jal.}{Becario de Desarrollo de Software}{Enero 2026 -- Abril 2026}
\begin{itemize}
  \item Participé en el desarrollo del sistema de control de peso para básculas industriales.
  \item Integré la comunicación entre app de escritorio, app móvil y Raspberry Pi.
  \item Desarrollo con Python, .NET, SQL Server y Android Studio; pruebas y documentación.
\end{itemize}

\entry{VirtuaMX}{Durango, México}{Becario en Desarrollo Web}{Mayo 2024 -- Agosto 2024}
\begin{itemize}
  \item Desarrollo de módulos frontend y backend; desarrollo y consumo de APIs REST.
  \item Git y GitHub bajo metodologías ágiles, cumpliendo entregables en tiempo y forma.
\end{itemize}

%==================== PROYECTOS ====================
\section{Proyectos destacados}
\begin{itemize}
  \item \textbf{MedicalRecords} (propio): sistema de expedientes médicos en \textbf{C\#/.NET} con SQL Server, con
        pruebas y documentación. Autor único.
  \item \textbf{Sistema de Asistencia} (PHP, MySQL): control de asistencia con roles, registro por QR y huella,
        esquema relacional y reportes exportables.
  \item \textbf{Sistema Ganadero Interactivo} (Python, Flask, SQLite): registro de ganado con exportación a Excel y PDF.
  \item \textbf{Scanner-QR} (React Native, Expo): app móvil de escaneo de QR con cámara, historial y almacenamiento local.
  \item \textbf{AgroBot -- HackDGO 2024} (equipo): robot agrícola con Arduino/ESP8266, Firebase y Google Vertex AI
        para análisis de imágenes de cultivos.
\end{itemize}

%==================== EDUCACIÓN ====================
\section{Educación}
\entry{Ingeniería en Desarrollo y Gestión de Software}{Durango}{Universidad Tecnológica de Rodeo}{2024 -- 2026}
\vspace{2pt}
\entry{TSU en Tecnologías de la Información -- Desarrollo de Software Multiplataforma}{Durango}{Universidad Tecnológica de Rodeo}{2022 -- 2024}

%==================== HABILIDADES ====================
\section{Habilidades técnicas}
\begin{itemize}
  \item \textbf{Backend:} C\#/.NET, ASP.NET Web API, Python, PHP, Java, Node.js/Express, Laravel, Flask, APIs REST, MVC.
  \item \textbf{Frontend:} JavaScript, TypeScript, Vue.js, React, HTML5, CSS3, Tailwind.
  \item \textbf{Bases de datos:} SQL Server, MySQL, PostgreSQL, MongoDB, SQLite; diseño relacional, respaldos y sincronización.
  \item \textbf{Móvil:} Java/Android, React Native, Expo, Android Studio.
  \item \textbf{Sistemas y herramientas:} Git, GitHub, Linux/Raspberry Pi, Windows Server, Postman, Scrum/Kanban.
  \item \textbf{Otros:} Ciberseguridad (fundamentos), ISO 27001, IoT (Arduino/ESP8266), Google Vertex AI.
\end{itemize}

%==================== CERTIFICACIONES ====================
\section{Certificaciones}
{\small
\begin{itemize}
  \item SQL y Bases de Datos desde Cero --- Udemy (2026)
  \item Artificial Intelligence Fundamentals --- IBM (2026)
  \item Cybersecurity Fundamentals --- IBM (2025)
  \item Internet de las Cosas (IoT) --- Santander Open Academy (2025)
  \item Seguridad digital para tu día a día --- Santander Open Academy (2025)
  \item Microsoft Copilot --- Santander Open Academy (2025)
  \item Introducción a la Ciencia de Datos --- Santander Open Academy (2024)
  \item IA y Productividad (Google) --- Santander Open Academy (2024)
  \item Python Essentials 1 y 2 --- Cisco (2023)
  \item EF SET English --- B2 (59/100), EF SET (2025)
  \item Liderazgo --- Capacítate para el Empleo (2026) \quad\textbullet\quad Pensamiento estratégico --- Santander Open Academy (2025)
  \item Marca personal 360\textdegree{} --- Santander Open Academy (2025) \quad\textbullet\quad Derechos Humanos y Criterios ESG --- ANUIES (2023)
\end{itemize}}

%==================== IDIOMAS ====================
\section{Idiomas}
Español (nativo) \quad\textbullet\quad Inglés --- B2 (Intermedio alto)

\end{document}

%==================================================================
% NOTA:
% - Compila directo con pdflatex u Overleaf (sin paquetes especiales).
% - Paleta cobre/tinta a juego con el portafolio.
%==================================================================
